\documentclass[11pt]{article}
\newcommand{\LabNumber}{11}
\newcommand{\LabTitle}{Merge Sort and Quick Sort}
\newcommand{\LabTopic}{Divide-and-conquer sorting, partitioning, complexity analysis}
\input{common.tex}

\begin{document}
\labheader

\section*{Objectives}
\begin{itemize}[leftmargin=*,nosep]
  \item Implement merge sort and understand its $O(n \log n)$ divide-and-conquer structure.
  \item Implement quick sort with the Lomuto partition scheme; analyse best, average, and worst cases.
  \item Measure empirical running times and compare with theoretical predictions.
  \item Apply sorting to practical problems: finding duplicates and the $k$-th smallest element.
\end{itemize}

\begin{summary}[Quick Reference]
\textbf{Merge sort.}
\begin{lstlisting}
void mergeSort(int[] a, int lo, int hi) {
    if (lo >= hi) return;
    int mid = (lo + hi) / 2;
    mergeSort(a, lo, mid);
    mergeSort(a, mid+1, hi);
    merge(a, lo, mid, hi);
}
\end{lstlisting}

\textbf{Quick sort pivot (Lomuto).}\quad Partition around \texttt{a[hi]}; elements $<$ pivot go left.\\
\textbf{Worst case.}\quad Already-sorted input with last-element pivot $\Rightarrow O(n^2)$.\\
\textbf{Fix.}\quad Swap a random element with \texttt{a[hi]} before partitioning.

\bigskip
\begin{tabular}{@{}lll@{}}
\toprule
Algorithm & Worst & Average \\
\midrule
Merge sort & $O(n \log n)$ & $O(n \log n)$ \\
Quick sort & $O(n^2)$ & $O(n \log n)$ \\
\bottomrule
\end{tabular}
\end{summary}

\subsection*{Quick Experiments (Try It)}

\begin{tryit}{Merge sort on paper first}
Trace merge sort on $[5, 2, 8, 1, 9, 3]$ by hand.
Draw the recursive split tree and the merge steps.
How many total element comparisons are made?
\end{tryit}

\begin{tryit}{Quick sort worst case}
Instrument your partition function with a comparison counter.
Compare the total comparisons for a sorted array of 10 vs.\ a random array of 10.
\end{tryit}

\section*{In-Lab Exercises (target: 2 hours)}

\begin{labexercise}{Implement Merge Sort \hfill {\small \itshape (40 min)}}
Implement \texttt{void mergeSort(int[] a, int lo, int hi)} and the helper
\texttt{void merge(int[] a, int lo, int mid, int hi)}.
\begin{itemize}[leftmargin=*,nosep]
  \item Use a temporary array of size \texttt{hi - lo + 1} inside \texttt{merge}.
  \item Print the array state after each top-level merge call.
  \item Test on: \texttt{\{5,2,8,1,9,3\}}, a reversed array, an already-sorted array.
\end{itemize}
\end{labexercise}

\begin{labexercise}{Implement Quick Sort \hfill {\small \itshape (40 min)}}
Implement \texttt{void quickSort(int[] a, int lo, int hi)} using the Lomuto partition
(pivot = \texttt{a[hi]}), where \texttt{int partition(int[] a, int lo, int hi)} returns
the final pivot index.
\begin{itemize}[leftmargin=*,nosep]
  \item Print the array after each \texttt{partition} call.
  \item Add a randomised-pivot variant.
  \item Compare call counts with merge sort on the same arrays.
\end{itemize}
\end{labexercise}

\begin{labexercise}{Empirical Comparison \hfill {\small \itshape (20 min)}}
Generate arrays of size 10\,000: random, sorted, reverse-sorted.
Time merge sort and quick sort (random pivot) using \texttt{System.nanoTime()}.
Record in a table. Does the ratio $T_{\text{merge}}/T_{\text{quick}}$ stay roughly constant?
\end{labexercise}

\begin{labexercise}{Find Duplicates \hfill {\small \itshape (20 min)}}
Given an unsorted \texttt{int[]}, find all duplicate values by sorting then scanning linearly.
Overall complexity? Compare with brute-force $O(n^2)$.
Test with \texttt{\{4, 1, 3, 2, 4, 7, 3\}} --- expected duplicates: \texttt{3, 4}.
\end{labexercise}

\begin{aicollab}
\textbf{Suggested prompts:}
\begin{itemize}[leftmargin=*,nosep]
  \item ``Trace the Lomuto partition on $[3,6,8,10,1,2,1]$ with pivot = 1.''
  \item ``Why does randomising the pivot make $O(n^2)$ behaviour negligibly unlikely?''
\end{itemize}
\medskip\textbf{Verify:} Ask the AI for the sorted result and pivot indices for a specific input.
Run your code on the same input and compare.
\end{aicollab}

\takehomebanner

\begin{trace}[Merge sort trace]
Trace \texttt{mergeSort(a, 0, 4)} on $a = [7, 3, 5, 1, 4]$.
Show the recursive call tree and array contents after each merge.
\end{trace}

\begin{trace}[Partition trace]
Trace \texttt{partition(a, 0, 5)} (Lomuto) on $a = [8, 3, 1, 5, 6, 2]$, pivot = \texttt{a[5] = 2}.
Show \texttt{i}, \texttt{j}, and the array after each iteration.
What is the final pivot index?
\end{trace}

\begin{writecode}[QuickSelect: $k$-th smallest]
Using quick sort's partition, write \texttt{int kthSmallest(int[] a, int k)}.
After partitioning at index $p$: if $p = k-1$ return \texttt{a[p]};
if $p > k-1$ recurse left; else recurse right.
This is \emph{QuickSelect}: $O(n)$ average.
\end{writecode}

\begin{drawex}{Merge sort recursion tree}
Draw the full recursion tree for merge sort on an array of 8 elements.
Label each node with its subarray range.
At each level, write the total comparisons in the merge step.
Show why the total is $O(n \log n)$.
\end{drawex}

\vspace{4pt}
\noindent\textbf{Submission.} Save files as \texttt{lab11\_ex*.java} and submit on the course site.

\end{document}
